%! Author = dave
%! Date = 7/16/24

% Preamble
\documentclass[11pt]{article}

% Packages
\usepackage{amsmath}

% Document
\begin{document}

    \section{Create Object}
    To use the history feature of the building tool, the object is associated with a text log.
    To create an object, use the create object button.
    This will place you into edit mode.
    NOTE: Change the mesh selection mode (top left of screen) from vertex to face before beginning operations.
    Otherwise, the redo-last feature doesn't work because the tool can't reset which face is selected.
    The first operation after creation is usually to insert a base polygon as the floor of the first room.

    \section{QArch Selection Mode}
    At the top of the QArch panel is a dropdown for selection mode.
    \begin{itemize}
        \item Single, each face selected gets a separate copy of the operation applied (so you can redo independently)
        \item Group, all faces selected get the sampe copy of the operation applies (redo one makes them all match)
        \item Region, some operations like buid-roof require the collection of faces to be treated as a one polygon region
    \end{itemize}

    \section{Coordinate System}
    Each polygon (face) has a coordinate system associated with it.
    For non-horizontal surfaces, the cross product between +Z and the face normal defines the +X direction.
    In other words, if you tip a square up on edge, the bottom edge becomes the X axis.
    For a horizontal surface, the global X axis becomes the face X axis.
    However, there are times when a different system makes sense and is imposed.
    When you create a frame (by making a smaller polygon inset into a larger polygon),
    the X axis follows the perimeter of the control polygon counter-clockwise around.
    This gives you a consistent way to reproduce scripts on any of the frame faces
    (since the key direction is usually in vs out).
    If you assign BT\_Nothing as the material, the X axis edge is colored blue, and the Y axis edge is colored red.

    \section{Size and Position}
    Many operations allow you to specify relative size or position.
    In this case, a 1 means the same size as the control.
    A special feature is that a negative size means to subtract that number from the control polygon size.
    Position is normally calculated from the bottom left corner of the polygon bounding box,
    but You may also have the option to calculate position relative to the center of the bounding box.

    \section{Deleting Faces}
    During construction, faces are never really deleted because they are used as references when executing the history
    script.
    Instead, the faces are tagged as DELETE, and marked hidden in the mesh edit mode.
    They also get a transparent material assigned.
    In object mode, all faces will be shown, and even though transparent, the existence of two faces in the same plane
    can cause rendering artifacts.
    The cleanup button will perform a true delete on the faces tagged DELETE .
    To restore the faces for further editing, simply use the rebuild button.

    \section{Set Active and Redo}
    When a face is selected, you can see the sequence of operations leading to the creation of that face with
    Set Active Operation or Redo Operation.
    Set Active Operation will highlight the faces created by the selected operation.
    Once highlighted, the Remove Operation button becomes active.
    All operations in the list following the removed operation are removed as well.
    You can select None as the active operation to deselect.
    Using the Redo Operation button brings up the redo-last panel in the lower left of the screen, just like when
    you first executed the operation.
    Some parameters may become locked if you have subsequent operations that depend on faces created by the operation
    you are redo-ing.
    Set Active Operation is also used before exporting a script to the catalog.
    The selected operation and all following operations become a script which you can apply to other locations.

    \section{Calc UV}
    Normally, UV coordinates are calculated for you, but if for some reason they are not, use this button to update
    the selected faces.

    \section{Cleanup}
    As noted in the section about deleted faces, the cleanup operation removes the hidden faces to enable proper
    rendering.

    \section{Insert Polygon}
    The insert polygon operation allows you to define a new shape in the coordinate system of a control shape, or the
    global coordinate system if no face is selected. Four options for joining the new geometery are given.
    \begin{itemize}
        \item Free, which simply positions a new polygon
        \item Inside, which clips the new polygon to only keep what is inside the boundary of the control polygon
        \item Outside, which clips the new polygon to keep the part outside the control
        \item Bridge, which joins the inside and outside polygons together and deletes the control polygon, assuming
        that the bridge polygons were meant to replace it.
    \end{itemize}
    The first point generated lies on the x axis, and the following points move counter-clockwise.
    The start-angle parameter lets you rotate the polygon.
    Set it to -180/n (if you have n sides) to get the first side parallel to the y axis and centered on the x axis.

    Other shapes besides regular polygons are possible.
    \begin{itemize}
        \item Self, make a scaled copy of the control polygon
        \item Ngon, make a regular polygon, optionally with an inset frame (keep the frame width smaller than the polygon!)
        \item Arch, make an arch, with or without built in frame
        \item Super, look up super-ellipse, this function is too complicated to describe and probably not useful
        \item Curve, choose a curve from the blend file
        \item Catalog, choose a curve from the qarch catalog
    \end{itemize}

    You may also provide an extrude distance which will float the new polygon over the control.

    \section{Plan Tools}
    These tools are under development for purposes of being able to construct a floorplan and then build it all at once
    in a selected style.

    \section{Macro Tools}
    Macro tools are really a wrapper around exported scripts that collect some user input to customize the script before
    adding it to the object history.

    \section{Apply Script}
    This button will load the catalog of scripts with descriptions and pictures.
    The script is applied as if you had executed all the individual steps.
    Sometimes, the script has a direction that makes sense on one side of a building but not the other.
    Use redo-operation to select the questionable step and customize it for your needs.

    \section{Import Mesh}
    Impor mesh lets you use an object in the current file, or a mesh from the catalog.
    It can be instanced as an array, or added directly to the mesh you are working on.
    When instanced, a geometry nodes modifier takes care of generating the instances.
    The input to the modifier uses a shared mesh so that if you edit one object, they all reflect the change.

    \section{Add Window}
    Add a window in small, medium, or large size.
    Customize fixed or sliding sash, number of panes, and shutters.
    A standard height is tried, but if the space is too small, the window will be centered vertically.

    \section{Niche}
    This is meant to be used in a brick wall, and the height and width are rounded to get brick sized results.
    Use to create a recess with an arched top.

    \section{Add Door}
    Doors come in one height but three widths.
    You can also customize whether it opens in or out, and which side the hinges are on.
    The door will be plain, but other scripts can be applied to the door face to decorate it.
    On the opposite side of the door, check the face orientation before applying a script.
    You may need to use flip-normal first.

    \section{Add Portico}
    This adds an overhang when applied to the space left above a door.
    The support columns are extruded from two small polygons.
    You can remove the last extrude operation and use import-mesh to add fancy columns instead.

    \section{Add Railing}
    A railing is made by extruding up a solid and then applying this function to the sides.
    When you are done, mark the unused sides as DELETE .
    You can use a simple shape extruded vertically, or select a curve to be rotated for each vertical spindle.

    \section{Build Roof}
    This uses a polygon skeleton algorithm to make complex hip roofs.
    It has a definite bug when working on L shaped roofing in some orientations.
    The roof tangent sets the angle of the roof, higher numbers are steeper.
    You may have to set the selection type to Single or Region for this.

    \section{Extend Gable}
    On a sloped face, usually a triangle generated by the build roof algorithm, this macro extrudes the face to create
    a vertical gable end.
    You may have to set the selection type to Single or Region for this.

    \section{Add Dormer}
    This will cut out a rectangle plus triangle from a sloped roof and extrude them to form a dormer.

    \section{Build Stairs}
    Stairs have an unusually large set of paramters.
    It is possible to make spiral stairs as well as straight stairs.
    Remember that z offset is away from the wall face, not up and down!

    \section{Detail Tools}
    The detail tools are low level routines that can be used to do almost anything.
    Be aware that there may be more than one way to do things, but probably only one of them will be a good way if you
    intend to export the script later.
    Especially avoid selecting faces from multiple operations when building scripts as the playback can get confused as
    to what faces already exist.
    If you are not saving a script, do whatever feels right.

    \section{Solidify Edges}
    A very powerful utility, this is used to create the mullions around window panes, the frames for doors, and the
    vertical posts in railings.
    You can specify which edge to solidify, or leave blank for all.
    A simple shape or a curve loaded from a catalog can be used.
    In fact, there are curves for Mansard roofing that can be applied to the attic floor to create a Mansard roof that
    wraps around with a nice profile.
    It is also possible to specify dash width and spacing.
    The dash feature makes dentil molding easy, and is used in some of the raised brick scripts.

    \section{Grid Divide}
    This feature can have x or y cuts set to zero so that you can split horizontally or vertically.
    Without defining the size, you get evenly spaced cuts.
    With a defined size, you set the spacing and get the number specified.
    This is how the door and window scripts set their position: they know how wide the frame has to be and set the cut
    number to 2, then let the user specify the offset from the center of the face.
    This makes a nice vertical slice of the right size.

    \section{Split Face}
    With split face you can pick two vertex indices to connect, or pick one vertex and a direction (x or y). The result
    is two faces.

    \section{Quoin Divide}
    Quoins are the large blocks set into the corners of brick buildings.
    The Quoin Divide cuts like a big set of zipper teeth so that you can change the material to a contrast on one side
    vs the other.

    \section{Extrude Fancy}
    In addition to extruding the face, you can resize it and twist while extruding.
    You can select a direction other than normal, and if so, whether to align the end face with that direction.

    \section{Extrude Sweep}
    The sweep is tricky to get used to.
    If you set the origin to <0,0,0> the rotation will be around the bottom left corner of the bounding box.
    If you set the axis to <1,0,0> the axis will be the bottom edge of the bounding box.
    Moving the origin to y=1 and is\_relative\_y=True will position the origin at the top of the bounding box.
    Just try it out.

    \section{Project Face}
    You choose the index of the target face by trial and error since it's hard to know what order the face selection
    will be presented in.
    The faces other than the target project along their normals until they meet the plane of the target.
    Or, you can use bridge to join to faces together.
    The hip option is used to join with a right angle, like making molding on one exterior wall join to molding on the
    next wall.

    \section{Build Face}
    Switch to edge or vertex mode, make a selection, and use this button to create a face from the selection.
    Then switch back to face selection mode.

    \section{Make Louvers}
    Louvers are used in shutters or if you want blinds for a window.
    They can also be connected for making simple stairs or other things.

    \section{Perpendicular Face}
    Sometimes you need to switch face orientation to build something.
    This generates a polygon flipped out 90 degrees from the control.

    \section{Set Oriented Material}
    The wood grain texture uses extra data to control how the face is oriented relative to the wood grain.
    This button takes the long direction of the face(s) and orients the wood in that direction.

    \section{Flip Normal}
    In addition to flipping the normal after an extrude, this can be used with the toggle off to create a no-action.
    The no-action may be a useful parent for a set of operations that you want to export to the catlog of scripts.

    \section{Face Info}
    With one face selected, face info reports on the tag, material, and other properties of interest.
    With one or more faces selected, clicking the button allows you to edit that property.
    Also reported is the operation and face number within the operation, sometimes useful for tracking down how a face
    was generated.


\end{document}